% 合并可视化图章节
% 说明：
% 1. 本文件设计为被主文档 \input{} 进来使用。
% 2. 这里将原图 5 和图 6 合并为一个跨栏图 5，避免重复说明同一件事。
% 3. 这里使用 cuted 宏包提供的 strip 环境，将合并后的图作为非浮动跨栏内容直接放在页顶。

% 统一控制合并图的单张小图宽度，后续如果需要给参考文献留更多空间，只需要略微调小该值。
\newcommand{\appendixablimg}[1]{\includegraphics[width=0.144\textwidth]{#1}}

% 手动生成跨栏区内的图题，避免 figure* 浮动体把图推到下一页。
\newcommand{\appendixablcaption}[2]{%
    \refstepcounter{figure}\label{#1}%
    \vspace{0.1em}%
    \parbox{0.98\textwidth}{\footnotesize 图~\thefigure. #2}%
}

\begin{strip}
\begin{center}
\vspace{-4.2em}
{\scriptsize
\setlength{\tabcolsep}{1pt}
\renewcommand{\arraystretch}{0.82}
\begin{tabular}{@{}c@{\hspace{2pt}}ccc@{\hspace{5pt}}ccc@{}}
& \multicolumn{3}{c}{Baseline} & \multicolumn{3}{c}{Ours} \\
\texttt{frame} & 0 次 & 4 次 & 8 次 & 0 次 & 4 次 & 8 次 \\
\texttt{0004} &
\appendixablimg{../../../images/self_compare/baseline/train_0004/times_0000.jpg} &
\appendixablimg{../../../images/self_compare/baseline/train_0004/times_0004.jpg} &
\appendixablimg{../../../images/self_compare/baseline/train_0004/times_0008.jpg} &
\appendixablimg{../../../images/self_compare/ours/train_0004/times_0000.jpg} &
\appendixablimg{../../../images/self_compare/ours/train_0004/times_0004.jpg} &
\appendixablimg{../../../images/self_compare/ours/train_0004/times_0008.jpg} \\
\texttt{0084} &
\appendixablimg{../../../images/self_compare/baseline/train_0084/times_0000.jpg} &
\appendixablimg{../../../images/self_compare/baseline/train_0084/times_0004.jpg} &
\appendixablimg{../../../images/self_compare/baseline/train_0084/times_0008.jpg} &
\appendixablimg{../../../images/self_compare/ours/train_0084/times_0000.jpg} &
\appendixablimg{../../../images/self_compare/ours/train_0084/times_0004.jpg} &
\appendixablimg{../../../images/self_compare/ours/train_0084/times_0008.jpg} \\
\texttt{0609} &
\appendixablimg{../../../images/self_compare/baseline/train_0609/times_0000.jpg} &
\appendixablimg{../../../images/self_compare/baseline/train_0609/times_0004.jpg} &
\appendixablimg{../../../images/self_compare/baseline/train_0609/times_0008.jpg} &
\appendixablimg{../../../images/self_compare/ours/train_0609/times_0000.jpg} &
\appendixablimg{../../../images/self_compare/ours/train_0609/times_0004.jpg} &
\appendixablimg{../../../images/self_compare/ours/train_0609/times_0008.jpg} \\
\texttt{0784} &
\appendixablimg{../../../images/self_compare/baseline/train_0784/times_0000.jpg} &
\appendixablimg{../../../images/self_compare/baseline/train_0784/times_0004.jpg} &
\appendixablimg{../../../images/self_compare/baseline/train_0784/times_0008.jpg} &
\appendixablimg{../../../images/self_compare/ours/train_0784/times_0000.jpg} &
\appendixablimg{../../../images/self_compare/ours/train_0784/times_0004.jpg} &
\appendixablimg{../../../images/self_compare/ours/train_0784/times_0008.jpg} \\
\end{tabular}
}
\par
\appendixablcaption{fig:ablation_visual}{前馈回归器消融的可视化结果。展示代表帧 \texttt{train\_0004}、\texttt{train\_0084}、\texttt{train\_0609} 与 \texttt{train\_0784} 在 0、4、8 次对当前帧的优化后的结果；每行左三列为 Baseline，右三列为 Ours。}
\vspace{0.4em}
\end{center}
\end{strip}

% 参考文献标题先占位，当前版本不放置任何实际引用条目。
\section*{References}
